\documentclass{article}
\usepackage{/home/user1/code/tex/sty}
\usepackage{amsfonts}
\usepackage{amsmath}
\usepackage{geometry}
\usepackage{graphicx}
\usepackage{hyperref}
\usepackage{floatrow}
\newcommand{\fig}[3]{
\begin{figure}[H]
\begin{center}
\includegraphics[height=#1]{#2}
\caption{#3}
\label{#2}
\end{center}
\end{figure}}
\usepackage[backend=bibtex]{biblatex}
\addbibresource{bib.bib}
\setlength{\parindent}{0pt}
\geometry{top=1in,bottom=1in,left=1in,right=1in}
\begin{document}
\begin{center}
\textbf{\Large{Problem Set 3 Correction}}\\~\\
\begin{tabular}{l l}
Student&Agustin Romero\\
Email&ar1@stanford.edu\\
Instructor&Caterina Vernieri\\
Course&Physics 252\\
Institution&Stanford University\\
Date&\today\\
\end{tabular}
\end{center}
\section*{Problem 1.a.}
The student solution was correct.
\section*{Problem 1.b.}
The student solution was not correct. Mesons have parity
\e{P=(-1)^{L+1}}
For $B^+$, $J^{P}=0^-$, so $B^+$ has odd parity. Change 
\e{\hat{P}\ket{B^+(\vec{p})}=\ket{B^+(-\vec{p})}}
to
\e{\boxed{\hat{P}\ket{B^+(\vec{p})}=-\ket{B^+(-\vec{p})}}}
\section*{Problem 1.c.}
The student solution was not correct. Photons have $I(J^{PC})=0,1(1^{--})$, so odd. Change
\e{\hat{P}\ket{\gamma(\vec{p})}=\ket{\gamma(-\vec{p})}}
to
\e{\boxed{\hat{P}\ket{\gamma(\vec{p})}=-\ket{\gamma(-\vec{p})}}}
\section*{Problem 1.d.}
The student solution was correct.
\section*{Problem 1.e.}
The student solution was correct.
\section*{Problem 2.a.}
The student solution was correct.
\section*{Problem 2.b.}
The student solution was correct.
\section*{Problem 2.c.}
The student solution was correct.
\section*{Problem 2.d.}
The student solution suggested that vector meson decay violates angular momentum conservation, but this is not true. $K^{0*}$ and $\overline{K^{0*}}$ can mix similarly to $K^{0}$ and $\overline{K^{0}}$, just at $J=1$ instead of $J=1$, but this is not observed, because \fbox{the $K^{0*}$ decays before mixing can occur.}
\section*{Problem 3.}
The student solution was correct.
\section*{Problem 4.a.}
The student solution was correct.
\section*{Problem 4.b.}
The student solution was correct.
\section*{Problem 4.c.}
The student solution was not correct. The student suggested that the spins of the two photons must be oriented oppositely such that the total angular momentum of the $\ket{\gamma \gamma}$ final state is 0, but this is not necessarily true. The $\eta$ has $I^G(J^{PC})=O^+(0^{-+})$, so $s=0$ and $l=1$. The two photon final state has $s=1$ and $l=1$ as shown in 4.b. The total angular momentum conservation requires that
\e{|l-s|\leq j \leq l+s}
So conservation of total angular momentum is \fbox{satisfied by the process $\eta\rightarrow \gamma \gamma$}.
\section*{Problem 5.a.}
The student solution was correct.
\section*{Problem 5.b.}
The student solution was correct.
\end{document}
